\documentclass[11pt]{article}

\usepackage{amsmath,amssymb}
\usepackage{xurl}
\usepackage[hidelinks]{hyperref}
\setlength{\emergencystretch}{2em}

% Shared commands for notes in docs/paper-gaps/.
% Notation follows the LDT paper (references/ldt-paper/).

\newcommand{\Lean}{\textsc{Lean}}
\newcommand{\leanid}[1]{\textsf{#1}}
\newcommand{\ghissue}[1]{\href{https://github.com/LionSR/MIPStarRE/issues/#1}{GitHub issue \##1}}
\newcommand{\ghpr}[1]{\href{https://github.com/LionSR/MIPStarRE/pull/#1}{GitHub PR \##1}}

% --- Paper-compatible notation ---
\newcommand{\F}{\mathbb{F}}
\newcommand{\N}{\mathbb{N}}
\newcommand{\C}{\mathbb{C}}
\newcommand{\tr}{\operatorname{tr}}
\newcommand{\ot}{\otimes}
\newcommand{\E}{\mathop{\bf E\/}}
\newcommand{\bx}{\mathbf{x}}
\newcommand{\by}{\mathbf{y}}
\newcommand{\bu}{\mathbf{u}}
\newcommand{\bv}{\mathbf{v}}

% Approximation relation (paper uses \approx_\eps and \simeq_\zeta)
\newcommand{\approxeps}{\approx_{\varepsilon}}
\newcommand{\simgeq}{\simeq_{\zeta}}

% Polynomial spaces (paper uses \polyfunc, \polysub, \polymeas)
\newcommand{\polyfunc}[3]{\operatorname{Poly}(#1,#2,#3)}
\newcommand{\polysub}[3]{\operatorname{PolySub}(#1,#2,#3)}
\newcommand{\polymeas}[3]{\operatorname{PolyMeas}(#1,#2,#3)}


\title{Issue 1099: The Sharper Local Repair for the Line-169 Transport Step}
\author{}
\date{2026-05-04}

\begin{document}
\maketitle

\section{Purpose}

This note records the sharper local correction to the line-169 transport step in
the main-induction projectivization chain of
\cite{Ji2020LowIndividualDegree}.\footnote{Tracked by \ghissue{1099}. The
surrounding projectivization argument is in
\path{references/ldt-paper/inductive_step.tex}, lines~130--185. The specific
line-169 replacement step is lines~167--173.}
The broader discrepancy note
\path{issue-1099-line169-triangle-sub-loss.tex} explains why the printed paper
argument cannot yield the exact error $\zeta_1$. The present note isolates the
sharper local repair that is actually used by the current formal base-case
assembly.

The point of this repair is that the generic corrected route
\[
  \zeta_1 + \sqrt{\zeta_2}
\]
is mathematically valid but quantitatively wasteful, because $\zeta_2$ is the
\emph{post-completion} error.  If one compares with the orthonormalized
projective submeasurement \emph{before} completion, then the square root lands
on the smaller orthonormalization error instead, producing the sharper bound
\[
  \zeta_1 + 10\zeta_1^{1/8}.
\]

\section{The setup}

At this stage of the induction, the argument has already produced two
polynomial-outcome measurements $G^{\mathrm A}$ and $G^{\mathrm B}$ satisfying
the line-130 consistency estimate
\[
  G^{\mathrm A}_g \otimes I \simeq_{\zeta_1} I \otimes G^{\mathrm B}_g.
  \tag{\path{eq:G-self-consistency}}
\]
Orthogonalization then supplies projective submeasurements
$P^{\mathrm A}$ and $P^{\mathrm B}$ with
\[
  G^{\mathrm A}_g \otimes I
  \approx_{100\zeta_1^{1/4}}
  P^{\mathrm A}_g \otimes I,
\]
and similarly on Bob's side.

Write
\[
  \epsilon := 100\zeta_1^{1/4}.
\]
Then the orthonormalization estimate is simply
\[
  G^{\mathrm A}_g \otimes I \approx_{\epsilon} P^{\mathrm A}_g \otimes I.
\]
Afterwards one completes $P^{\mathrm A}$ to a projective measurement
$Q^{\mathrm A}$ by adding the missing mass $I-\sum_g P^{\mathrm A}_g$ at a
distinguished outcome.

\section{The sharp local correction}

The correction proceeds in two steps.

\subsection*{Step 1: compare before completion}

Fix the partner measurement $G^{\mathrm B}$.  Proposition
\path{prop:easy-approx-from-approx-delta} bounds the change in diagonal match
mass when one replaces $G^{\mathrm A}$ by a $\approx_{\epsilon}$-close family.
Specialized to the constant single-question setting, it gives
\[
  \left|
    \sum_g \langle \psi,
      (G^{\mathrm A}_g \otimes G^{\mathrm B}_g)\psi\rangle
    -
    \sum_g \langle \psi,
      (P^{\mathrm A}_g \otimes G^{\mathrm B}_g)\psi\rangle
  \right|
  \le \sqrt{\epsilon}.
\]
Since line~130 says that the first match mass is at least $1-\zeta_1$, one
obtains
\[
  \sum_g \langle \psi,
    (P^{\mathrm A}_g \otimes G^{\mathrm B}_g)\psi\rangle
  \ge 1 - \zeta_1 - \sqrt{\epsilon}.
\]
Equivalently,
\[
  P^{\mathrm A}_g \otimes I
  \simeq_{\zeta_1 + \sqrt{\epsilon}}
  I \otimes G^{\mathrm B}_g.
\]

\subsection*{Step 2: complete without further loss}

Now pass from $P^{\mathrm A}$ to the completed projective measurement
$Q^{\mathrm A}$.  The completion changes exactly one outcome by adding the
positive operator
\[
  I - \sum_g P^{\mathrm A}_g.
\]
Therefore the diagonal match mass against the fixed partner measurement
$G^{\mathrm B}$ can only increase.  No new negative term appears in the
consistency defect.  Hence
\[
  Q^{\mathrm A}_g \otimes I
  \simeq_{\zeta_1 + \sqrt{\epsilon}}
  I \otimes G^{\mathrm B}_g.
\]

Substituting the value of $\epsilon$ gives the explicit form
\[
  Q^{\mathrm A}_g \otimes I
  \simeq_{\zeta_1 + \sqrt{100\zeta_1^{1/4}}}
  I \otimes G^{\mathrm B}_g,
\]
that is,
\[
  Q^{\mathrm A}_g \otimes I
  \simeq_{\zeta_1 + 10\zeta_1^{1/8}}
  I \otimes G^{\mathrm B}_g.
  \tag{sharper repaired line 169}
\]

The Bob-side mirror is identical:
\[
  I \otimes Q^{\mathrm B}_g
  \simeq_{\zeta_1 + 10\zeta_1^{1/8}}
  G^{\mathrm A}_g \otimes I.
\]

\section{Why this is sharper than the generic correction}

The generic corrected route applies \path{prop:triangle-sub} only after the
completion estimate
\[
  G^{\mathrm A}_g \otimes I \approx_{\zeta_2} Q^{\mathrm A}_g \otimes I,
\]
and therefore yields
\[
  Q^{\mathrm A}_g \otimes I
  \simeq_{\zeta_1 + \sqrt{\zeta_2}}
  I \otimes G^{\mathrm B}_g.
\]
This is correct, but $\zeta_2$ is already the larger post-completion scalar.  In
contrast, the local repair takes the square root earlier, at the smaller
orthonormalization scale $\epsilon=100\zeta_1^{1/4}$.  Thus
\[
  \zeta_1 + 10\zeta_1^{1/8}
\]
is strictly better than
\[
  \zeta_1 + \sqrt{\zeta_2}
\]
for the present cascade.

This quantitative distinction is exactly why the sharper repair preserves the
original final envelope used by the current formal development, whereas the
generic corrected route would force the weaker $1/65536$ scale.

\section{Relation with the formal development}

The formal development proves this sharper local repair directly.\footnote{The
general pre-completion transport statements are
\leanid{MIPStarRE.LDT.MakingMeasurementsProjective.ProjectivizationLine169Repair.leftConsistency\_of\_completion\_and\_sdd}
and
\leanid{MIPStarRE.LDT.MakingMeasurementsProjective.ProjectivizationLine169Repair.rightConsistency\_of\_completion\_and\_sdd}
in
\doclink{MIPStarRE/LDT/MakingMeasurementsProjective/ProjectivizationChain.lean}.
Specializing the error to
\leanid{MIPStarRE.LDT.MakingMeasurementsProjective.orthonormalizationError}
gives
\leanid{MIPStarRE.LDT.MakingMeasurementsProjective.ProjectivizationLine169Repair.leftConsistency\_with\_orthonormalization\_loss}
and its Bob-side mirror.  The corresponding Step-8 absorption is
\leanid{MIPStarRE.LDT.Test.repairedLine169PointError\_le\_mainFormalError} in
\doclink{MIPStarRE/LDT/Test/MainTheorem.lean}.}

In the current code, the active base-case target is the match-mass bridge
\leanid{MIPStarRE.LDT.Test.MainFormalBaseBranchCompletionObligations}, and the final
transport is factored through
\leanid{MIPStarRE.LDT.Test.mainFormal\_ofProjectiveCompletionResidual}.  The
older repaired base-case assembly theorem has been removed from the live
MainFormal route.  The generic corrected route
$\zeta_1 + \sqrt{\zeta_2}$ remains mathematically available, but it is not the
active route for the present formal theorem.

\section{Conclusion}

The sharper local correction is obtained by inserting the square root at the
orthonormalization stage rather than at the completion stage.  Writing
\(
\epsilon = 100\zeta_1^{1/4}
\), the corrected line-169 transport is
\[
  Q^{\mathrm A}_g \otimes I
  \simeq_{\zeta_1 + \sqrt{\epsilon}}
  I \otimes G^{\mathrm B}_g,
\]
equivalently,
\[
  Q^{\mathrm A}_g \otimes I
  \simeq_{\zeta_1 + 10\zeta_1^{1/8}}
  I \otimes G^{\mathrm B}_g.
\]

This does not recover the paper's printed exact $\zeta_1$ line, but it is a
strictly sharper correction than the generic route
$\zeta_1 + \sqrt{\zeta_2}$ and is sufficient to preserve the original final
error envelope in the current formal base-case assembly.

\bibliographystyle{alpha}
\bibliography{references}

\end{document}
